\documentclass[a4paper,11pt]{article}%

\usepackage{fullpage}%
\usepackage{url}%
\usepackage[T1]{fontenc}%
\usepackage[utf8]{inputenc}%
\usepackage[main=french]{babel}% % Adjust the main language

\usepackage{graphicx}%

\title{Activité "Puzzle Mathématique"}%
\author{Emile Thuiller \and Edouard Aubert}%
\date{}%

\begin{document}%

\maketitle%

\section{Résumé de l'activité}

\paragraph{Objectif}
Traduire des images en suite de nombres et inversement

\url{https://people.irisa.fr/Martin.Quinson/Mediation/InfoSansOrdi/} (voir jour 16)

\paragraph{Compétence}
Résoudre des puzzles et des problèmes difficiles

\paragraph{Niveau}
CE2 et plus


\paragraph{Durée}
Environ 45 minutes.

\section{Déroulement}

\begin{center}
\begin{tabular}{|l|p{5.5em}|p{15em}|p{6em}|p{6em}|}
  \hline
    \textbf{Durée} 
  & \textbf{Phases} 
  & \textbf{Activité et consigne} 
  & \textbf{Organisation} 
  & \textbf{Matériel} 
  \\
  \hline
  \hline
  2-3' 
  & Présentation de l'activité 
  & "Aujourd'hui nous allons résoudre des puzzles avec des pièces carrées"\newline
  & Oral Collectif 
  & (Aucun)
  \\
  \hline
  2-3'
  & Consignes 
  & Donner un exemple pour comprendre quand est-ce que deux pièces peuvent aller ensemble
  & Oral Collectif 
  & (Aucun)
  \\
  \hline
  18'
  & Mise en pratique
  & Vérifier dans chaque groupe que les consignes sont bien comprises 
  & En groupes de 3 à 5
  & 24 pièces carrées pour faire un puzzle
  \\
  \hline
  5'
  & Remise en commun
  & On demande si un groupe a trouvé une solution au puzzle. Ensuite, nouvelle consigne : on leur demande une méthode permettant de construire toutes les combinaisons possibles des pièces
  & Oral collectif
  & D'autres grilles vierges et crayons de couleurs
  \\
  \hline
  7'
  & Recherche
  & Guider les groupes les plus en retrait, réexpliquer les consignes au besoin chez les groupes désorientés.
  & En groupes de 3 à 5, comme à la mise en pratique
  & (Aucun)
  \\
  \hline
  5'
  & Restitution
  & Demander à la classe si elle a trouvé une solution
  & Oral collectif
  & (Aucun)
  \\
  \hline
  5'
  & Insitutionali-sation
  & Problèmes obligeant à faire une recherche exhaustive, retour sur trace, exemple avec l'écran de veille faisant des pavages avec polyominaux (via xscreensaver, utiliser xorg).
  & Oral collectif
  & (Aucun)
  \\
  \hline
\end{tabular}
\end{center}

\end{document}%
